\textbf{Phần 1. Trắc nghiệm lựa chọn}
\setcounterpageref{ex}{0}

\begin{ex} %[2D1V05]
Tính nhanh giá trị của biểu thức: $A = 24 + 15 + 36 + 25$.
\choice
{$90$}
{$95$}
{\True $100$}
{$105$}
\loigiai{
Gộp các số tạo số tròn chục: $A = (24 + 36) + (15 + 25) = 60 + 40 = 100$.
\textbf{Đáp án của câu hỏi là C}
}
\end{ex}

\begin{ex} %[2D1V05]
Tìm $x$ biết: $x + 23 = 78 - 15$.
\choice
{$35$}
{\True $40$}
{$45$}
{$50$}
\loigiai{
Vế phải: $78 - 15 = 63$. Ta có: $x + 23 = 63 \Rightarrow x = 63 - 23 = 40$.
\textbf{Đáp án của câu hỏi là B}
}
\end{ex}

\begin{ex} %[2D1V05]
Một xe buýt chở $36$ hành khách. Đến trạm dừng có $12$ người xuống xe và $8$ người lên xe. Hỏi trên xe buýt lúc này có bao nhiêu hành khách?
\choice
{$30$ người}
{$34$ người}
{$28$ người}
{\True $32$ người}
\loigiai{
Số người trên xe lúc này là: $36 - 12 + 8 = 24 + 8 = 32$ (người).
\textbf{Đáp án của câu hỏi là D}
}
\end{ex}

\begin{ex} %[2D1V05]
Điền số thích hợp vào ô trống: $96 - \dots = 42 + 21$.
\choice
{\True $33$}
{$35$}
{$31$}
{$23$}
\loigiai{
Vế phải: $42 + 21 = 63$. Ô trống $= 96 - 63 = 33$.
\textbf{Đáp án của câu hỏi là A}
}
\end{ex}

\begin{ex} %[2D1V05]
Tính hiệu của số lớn nhất có hai chữ số và tổng của hai số $32$ và $25$.
\choice
{$38$}
{$40$}
{\True $42$}
{$44$}
\loigiai{
Số lớn nhất có hai chữ số là 99. Tổng hai số là $32 + 25 = 57$.\nHiệu cần tìm: $99 - 57 = 42$.
\textbf{Đáp án của câu hỏi là C}
}
\end{ex}

\begin{ex} %[2D1V05]
Một cửa hàng có $85$ mét vải. Buổi sáng bán được $23$ mét, buổi chiều bán được $31$ mét. Hỏi cửa hàng còn lại bao nhiêu mét vải?
\choice
{$29\text{ m}$}
{\True $31\text{ m}$}
{$33\text{ m}$}
{$25\text{ m}$}
\loigiai{
Tổng số mét vải đã bán: $23 + 31 = 54\text{ (m)}$.\nSố vải còn lại: $85 - 54 = 31\text{ (m)}$.
\textbf{Đáp án của câu hỏi là B}
}
\end{ex}

\begin{ex} %[2D1V05]
Tìm một số biết rằng lấy số đó trừ đi $14$ rồi cộng thêm $23$ thì được $68$.
\choice
{$57$}
{$61$}
{$55$}
{\True $59$}
\loigiai{
Tính ngược từ cuối: $68 - 23 + 14 = 45 + 14 = 59$.
\textbf{Đáp án của câu hỏi là D}
}
\end{ex}

\begin{ex} %[2D1V05]
Cho phép tính: $\overline{a4} + 35 = 79$. Giá trị của chữ số $a$ là
\choice
{\True $4$}
{$3$}
{$5$}
{$2$}
\loigiai{
Ta có: $\overline{a4} = 79 - 35 = 44$. Do đó $a = 4$.
\textbf{Đáp án của câu hỏi là A}
}
\end{ex}

\begin{ex} %[2D1V05]
Điền dấu thích hợp vào chỗ chấm: $68 - 25 \dots 21 + 22$.
\choice
{$>$}
{$<$}
{\True $=$}
{Không so sánh được}
\loigiai{
Vế trái: $68 - 25 = 43$. Vế phải: $21 + 22 = 43$. Vì $43 = 43$ nên điền dấu $=$.
\textbf{Đáp án của câu hỏi là C}
}
\end{ex}

\begin{ex} %[2D1V05]
Bạn Linh có $58$ con tem. Linh cho Nga $15$ con tem và cho Mai $12$ con tem. Hỏi Linh còn lại bao nhiêu con tem?
\choice
{$33$ con tem}
{\True $31$ con tem}
{$29$ con tem}
{$35$ con tem}
\loigiai{
Linh cho hai bạn tất cả: $15 + 12 = 27$ (con tem).\nLinh còn lại: $58 - 27 = 31$ (con tem).
\textbf{Đáp án của câu hỏi là B}
}
\end{ex}

\begin{ex} %[2D1V05]
Có ba thùng sách. Thùng thứ nhất có $21$ quyển, thùng thứ hai có $24$ quyển, thùng thứ ba có số sách bằng tổng hai thùng đầu trừ đi $15$ quyển. Thùng thứ ba có bao nhiêu quyển sách?
\choice
{$35$ quyển}
{$28$ quyển}
{$32$ quyển}
{\True $30$ quyển}
\loigiai{
Tổng hai thùng đầu: $21 + 24 = 45$ (quyển).\nThùng thứ ba có: $45 - 15 = 30$ (quyển).
\textbf{Đáp án của câu hỏi là D}
}
\end{ex}

\begin{ex} %[2D1V05]
Tìm hiệu của số chẵn lớn nhất có hai chữ số và số lẻ nhỏ nhất có hai chữ số.
\choice
{\True $87$}
{$89$}
{$88$}
{$86$}
\loigiai{
Số chẵn lớn nhất có hai chữ số là 98. Số lẻ nhỏ nhất có hai chữ số là 11.\nHiệu là: $98 - 11 = 87$.
\textbf{Đáp án của câu hỏi là A}
}
\end{ex}

\textbf{Phần 2. Trắc nghiệm đúng sai}
\setcounterpageref{ex}{0}

\begin{ex} %[2D1V05]
Xét tính đúng sai của các biểu thức tính toán sau:
\choiceTF
{\True Biểu thức $45 + 23 - 12$ có giá trị bằng $56$.}
{Biểu thức $89 - 34 - 15$ có giá trị bằng giá trị của $89 - (34 - 15)$.}
{\True Tổng của số chẵn và số lẻ luôn luôn là số lẻ.}
{\True Trong phép tính $78 - 35 = 43$, các chữ số hàng đơn vị thỏa mãn $8 - 5 = 3$.}
\loigiai{
\begin{itemize}
    \item Ý a) Đúng: $45 + 23 = 68$; $68 - 12 = 56$.
    \item Ý b) Sai: $89 - 34 - 15 = 40$, còn $89 - 19 = 70$. Hai kết quả khác nhau.
    \item Ý c) Đúng: Chẵn + Lẻ = Lẻ.
    \item Ý d) Đúng: Hàng đơn vị: $8 - 5 = 3$. Đúng.
\end{itemize}
\begin{itemize}
    \item a) \textbf{Đúng}.
    \item b) \textbf{Sai}.
    \item c) \textbf{Đúng}.
    \item d) \textbf{Đúng}.
\end{itemize}
}
\end{ex}

\begin{ex} %[2D1V05]
Trong vườn nhà bác An có 45 cây ăn quả gồm cam, bưởi và xoài. Trong đó có 21 cây cam, số cây bưởi ít hơn cây cam 10 cây:
\choiceTF
{\True Số cây bưởi là 11 cây.}
{\True Tổng số cây cam và bưởi là 32 cây.}
{Số cây xoài trong vườn là 15 cây.}
{\True Số cây cam nhiều hơn số cây xoài là 8 cây.}
\loigiai{
\begin{itemize}
    \item Ý a) Đúng: $21 - 10 = 11$ cây bưởi.
    \item Ý b) Đúng: $21 + 11 = 32$ cây.
    \item Ý c) Sai: Số cây xoài là: $45 - 32 = 13$ cây (không phải 15 cây).
    \item Ý d) Đúng: $21 - 13 = 8$ cây.
\end{itemize}
\begin{itemize}
    \item a) \textbf{Đúng}.
    \item b) \textbf{Đúng}.
    \item c) \textbf{Sai}.
    \item d) \textbf{Đúng}.
\end{itemize}
}
\end{ex}

\textbf{Phần 3. Trả lời ngắn}
\setcounterpageref{ex}{0}

\begin{ex} %[2D1V05]
Tính nhanh giá trị của biểu thức: $12 + 25 + 38 + 15$.
\shortans{90}
\loigiai{
\textbf{Đáp số:} $90$\\
$(12 + 38) + (25 + 15) = 50 + 40 = 90$.
}
\end{ex}

\begin{ex} %[2D1V05]
Tìm x biết: $98 - x = 24 + 31$.
\shortans{43}
\loigiai{
\textbf{Đáp số:} $43$\\
$98 - x = 55 \Rightarrow x = 98 - 55 = 43$.
}
\end{ex}

\begin{ex} %[2D1V05]
Một mảnh đất hình chữ nhật có chu vi là 88 m. Nửa chu vi mảnh đất đó bằng bao nhiêu mét?
\shortans{44}
\loigiai{
\textbf{Đáp số:} $44$\\
Nửa chu vi: $88 : 2 = 44\text{ m}$.
}
\end{ex}

\begin{ex} %[2D1V05]
Nhà Mai thu hoạch được 56 kg cà chua, đã bán đi 34 kg. Số cà chua còn lại chia đều vào 2 rổ. Hỏi mỗi rổ có bao nhiêu ki-lô-gam cà chua?
\shortans{11}
\loigiai{
\textbf{Đáp số:} $11$\\
Số cà chua còn lại: $56 - 34 = 22\text{ kg}$. Mỗi rổ có: $22 : 2 = 11\text{ kg}$.
}
\end{ex}

\begin{ex} %[2D1V05]
Tìm một số biết rằng lấy số đó cộng với 32 rồi trừ đi 15 thì được kết quả là 64.
\shortans{47}
\loigiai{
\textbf{Đáp số:} $47$\\
Số đó là: $64 + 15 - 32 = 79 - 32 = 47$.
}
\end{ex}

\begin{ex} %[2D1V05]
Điền chữ số thích hợp thay cho dấu hỏi chấm: $?4 + 23 = 87$.
\shortans{6}
\loigiai{
\textbf{Đáp số:} $6$\\
Số đó là $87 - 23 = 64$. Vậy dấu ? là chữ số 6.
}
\end{ex}
